% Analysis section — deeper dives beyond main results

\subsection{Why Do Single-Path Controllers Fail?}
\label{sec:analysis:single_path}

A natural question is why simpler approaches---extracting difficulty signals from a single reasoning trace---consistently fail.
Table~\ref{tab:negative_results} summarizes our systematic evaluation of single-path controllers.

\begin{table}[h]
\centering
\caption{Systematic negative results for single-path difficulty controllers on GSM8K.
All methods use the same model (Qwen3.5-27B) and evaluation protocol.
$\Delta$ is accuracy change relative to the best fixed baseline.}
\label{tab:negative_results}
\begin{tabular}{lcc}
\toprule
Controller & Accuracy (\%) & $\Delta$ (pp) \\
\midrule
Best Fixed Budget (256 tokens) & XX.X & --- \\
\midrule
\textit{Feature-based:} \\
\quad Answer Presence + Token Util. + Consistency & XX.X & $-6.5$ \\
\quad Token Utilization Only & XX.X & $-4.2$ \\
\quad Answer Presence Only & XX.X & $-3.1$ \\
\midrule
\textit{Uncertainty-based:} \\
\quad Logit Entropy (requires white-box) & XX.X & $-6.4$ \\
\midrule
\textit{Search-based:} \\
\quad MCTS-R (serial explore-refine-verify) & XX.X & $-5.6$ \\
\midrule
\textbf{\method{} (ours, K=XX, probe=XXX)} & \textbf{XX.X} & \textbf{+X.X} \\
\bottomrule
\end{tabular}
\end{table}

The fundamental issue is \emph{information poverty}: a single reasoning trace provides at most $H(X_1)$ bits of information about problem difficulty, while $K$ independent traces provide up to $K \cdot H(X_1)$ bits (Proposition~\ref{prop:information}).
In practice, the information gain from single-path features is severely limited because:
(i)~answer presence is a noisy binary signal confounded by generation length;
(ii)~token utilization conflates difficulty with verbosity;
(iii)~logit-based uncertainty requires white-box access unavailable in API settings.

\subsection{Consensus Signal Quality}
\label{sec:analysis:consensus}

We validate the consensus signal against ground-truth difficulty using our Phase~0 analysis on existing multi-budget data (Section~\ref{sec:theory}).

\paragraph{Correlation with difficulty.}
The Spearman rank correlation between consensus score $C_K$ and problem difficulty is $\rho = -0.51$ ($p < 10^{-88}$), substantially stronger than any single-path feature (all $|\rho| < 0.2$).
This confirms Proposition~\ref{prop:information}: multi-path consensus is a fundamentally richer difficulty signal.

\paragraph{Consensus precision.}
When all $K$ paths agree on the same answer (full consensus), the answer is correct \textbf{72.8\%} of the time.
This high precision justifies the ``easy'' route: accepting the consensus answer without further computation.

\paragraph{Route calibration.}
Figure~\ref{fig:route_calibration} shows the accuracy of each route as a function of the consensus score $C_K$.
The easy route (high consensus) achieves XX.X\% accuracy, the medium route achieves XX.X\%, and the hard route achieves XX.X\%.
The monotonic relationship confirms that our threshold-based routing is well-calibrated.

\begin{figure}[h]
\centering
% PLACEHOLDER — will be generated from results
\fbox{\parbox{0.8\textwidth}{\centering\vspace{2cm}\textit{[Figure: Route calibration — accuracy vs consensus score, colored by route]}\vspace{2cm}}}
\caption{Accuracy as a function of consensus score. Higher consensus reliably indicates easier problems, validating the routing mechanism.}
\label{fig:route_calibration}
\end{figure}

\subsection{Token Efficiency Analysis}
\label{sec:analysis:efficiency}

\method{} achieves token savings primarily through the easy route, which avoids the costly resolution phase entirely.

\paragraph{Per-route token cost.}
\begin{itemize}[nosep]
\item \textbf{Easy} (XX\% of samples): $K \times B_\text{probe} = $ XX tokens. \textit{No resolution cost.}
\item \textbf{Medium} (XX\% of samples): $K \times B_\text{probe} + B_\text{ext} \approx $ XX tokens.
\item \textbf{Hard} (XX\% of samples): $K \times B_\text{probe} + B_\text{hard} = $ XX tokens.
\end{itemize}

\paragraph{Comparison to budget-matched baselines.}
A fair comparison accounts for the total tokens consumed.
\method{} uses an average of XX tokens per sample.
A fixed budget of XX tokens (the same average) achieves only XX.X\% accuracy, while \method{} achieves XX.X\%---a \textbf{+X.X\,pp} improvement at the same token budget.
This demonstrates that \method{}'s gains come from \emph{better allocation}, not simply \emph{more tokens}.

\paragraph{Wall-clock latency.}
Although \method{} generates $K$ exploration paths, these are \emph{embarrassingly parallel}: on hardware with sufficient memory bandwidth, the $K$ paths can be generated as a single batch, reducing wall-clock latency to approximately $B_\text{probe}$ tokens (not $K \times B_\text{probe}$).
In our sequential implementation, the latency overhead is $K\times$ for the exploration phase, but the resolution phase uses only a single generation call.

\subsection{Error Analysis}
\label{sec:analysis:errors}

We categorize errors by route to understand where \method{} fails:

\begin{itemize}[nosep]
\item \textbf{Easy-route errors} (false confidence): All paths agree on the \emph{wrong} answer.
This occurs when the problem has a common misconception that tricks the model consistently.
These errors are \emph{inherent}---no amount of additional compute would help, as the model lacks the capability.

\item \textbf{Medium-route errors} (insufficient extension): The extension phase fails to correct the initial partial answer.
Often caused by the extension continuing a flawed reasoning chain rather than starting fresh.

\item \textbf{Hard-route errors} (deliberation failure): Despite seeing all path summaries, the deliberation produces an incorrect answer.
These represent genuinely difficult problems where the model struggles even with full context.
\end{itemize}
